\begin{papercard}{wharton1994}
\cf{Stages} \stagetag{W}. The practitioner's guide to the \emph{third}, simplified version of the cognitive walkthrough (CW), the canonical source for the four questions and the ``credible story'' procedure. \emph{Note:} the appended copy is the chapter excerpt (printed pp.\ 105--140), and page numbers below are relative to the excerpt.

\cf{Purpose and scope} An inspection method evaluating ease of learning \emph{by exploration}, motivated by users preferring to learn while doing their tasks \pc{wharton1994}{1}{usability inspection method that focuses on evaluating}. It finds mismatches between users' and designers' task conceptions, poor labels and inadequate feedback, and it surfaces implicit assumptions about users' knowledge \pc{wharton1994}{3}{The method finds mismatches between users}. Narrow by design: it trades coverage of other usability attributes (consistency, catastrophic errors) for depth on learnability, so it should be combined with other methods \pc{wharton1994}{4}{These studies show that the method is narrowly}.

\cf{Process overview (Table 5.1)} (1) Define inputs: users, sample tasks, action sequences, interface description. (2) Convene analysts. (3) Walk through each action and tell a credible story using the four questions. (4) Record critical information: user knowledge requirements, assumptions about users, side issues and design changes, the success story. (5) Revise the interface \pc{wharton1994}{2}{Overview of the cognitive walkthrough process}. Two phases, preparatory and analysis \pc{wharton1994}{4}{The cognitive walkthrough analysis has two phases}.

\cf{Inputs (preparatory phase)} \emph{Users}: a description including background experience and knowledge of the task and the interface \pc{wharton1994}{5}{Who will be the users of the system}. \emph{Tasks}: a representative collection of benchmark tasks from core functionality plus combinations; concrete and realistic, with context such as database contents \pc{wharton1994}{6}{A critical question is how to select these representative tasks}. \emph{Action sequence}: how users view the task before learning the interface, plus the correct actions, at a granularity matched to user expertise (``same level as a successful prompt or effective tutorial'') \pc{wharton1994}{6}{What is the correct action sequence for each task and how is it described}. \emph{Interface}: the prompts before every action and the reactions after it \pc{wharton1994}{7}{The definition of the interface must describe the}.

\cf{Analysis: credible stories} For each step, analysts tell a credible story of why the expected user would choose the action, based on CE+ exploratory learning: rough task description, explore, observe the reaction, decide the next action \pc{wharton1994}{7}{The problem-solving process is described by Polson and Lewis}. Users often follow \emph{label-following} \pc{wharton1994}{8}{which leads them to select an action if the label for that action matches}. \textbf{The four questions}: Will the users try to achieve the right effect? Will the user notice that the correct action is available? Will the user associate the correct action with the effect? If the correct action is performed, will the user see that progress is being made? \pc{wharton1994}{8}{Will the users try to achieve the right effect}. These are guiding criteria, not absolute requirements \pc{wharton1994}{8}{These questions are loose guidelines}.

\cf{Recording} Group-visible displays: one for the group story, one per user class (what the user must know before and should learn during the task), one for side issues and design changes, plus video where possible \pc{wharton1994}{9}{What the user must know prior to performing the task}.

\cf{Success and failure stories} Success needs all four criteria; failure stories usually fail one, so failures are organised by criterion \pc{wharton1994}{12}{Success stories require success under all four of the analysts}. Ways users know things include experience, visible devices, labels, ``all other actions look wrong'', and recognisable system responses \pc{wharton1994}{11}{Users May Know What Effect To Achieve}. Also watch for time-outs, physically difficult actions and dropped terminator actions. Be suspicious of stories relying on system terminology or complex plans \pc{wharton1994}{14}{Some Other Problems to Watch For}. A worked phone-forwarding example shows step-by-step coding \pc{wharton1994}{14}{The task is to forward calls on a campus phone system}.

\cf{Staying on track} The analysis always follows the \emph{correct} path. Problems are noted, then the analysis proceeds as if the correct action had been taken, from the correct state, so what the user knows a priori or learns must be recorded clearly \pc{wharton1994}{19}{The previous example demonstrates that the analysis always tracks the correct}. Fixes are organised by failed criterion (eliminate the action or add a prompt; a more visible control; better labels; better feedback) \pc{wharton1994}{20}{Will the user be trying to achieve the right effect? If the interface fails on this}.

\cf{Evolution and evidence} Version 1 was a one-page form that failed with untrained analysts, who confused goals with actions. Version 2 was formal and detailed but too tedious. Version 3 uses credible stories and needs basic cognitive-theory knowledge, through training or a ``champion'' \pc{wharton1994}{23}{Our first version of the walkthrough used a single-page form}, \pc{wharton1994}{26}{believe that a basic understanding of the cognitive theory is essential}. No formal validation of version 3 yet \pc{wharton1994}{25}{We do not yet have any empirical data validating the new}. Comparative studies: CW is narrow and finds specific rather than general problems \pc{wharton1994}{30}{Because the walkthrough is}. It best predicted problems affecting real users in Cuomo's follow-up (58\,\%) \pc{wharton1994}{29}{walkthrough performed the best in terms of predicting problems which}. Group walkthroughs need a skilled facilitator, and task selection is critical \pc{wharton1994}{32}{noted, however, is the need for a good facilitator}.

\cf{Use for this thesis} Source text for the walkthrough protocol: define the user (e.g.\ a first-time German-speaking tracker user), privacy tasks (find the policy, understand sharing, opt out, delete, export), the correct action sequence, and the interface state. For every step, record the four answers with a credible success or failure story, plus knowledge requirements. Adapting CW from learnability to \emph{finding and understanding privacy information} must be stated as an adaptation.
\end{papercard}
